\documentclass[10pt,letterpaper]{article}
\usepackage[utf8]{inputenc}
\usepackage[T1]{fontenc}
\usepackage{lmodern}
\usepackage[letterpaper,landscape,top=1.5cm,bottom=1.5cm,left=1.5cm,right=1.5cm]{geometry}
\usepackage{amsmath}
\usepackage{booktabs}
\usepackage{tabularx}
\usepackage{xltabular}
\usepackage{xcolor}
\usepackage{fancyhdr}
\usepackage{lastpage}
\usepackage{caption}
\captionsetup[table]{name=Cuadro,labelfont=bf,textfont=bf,skip=4pt}

\definecolor{navyblue}{RGB}{26, 79, 124}
\pagestyle{fancy}
\fancyhf{}
\lhead{\small\color{navyblue}\textbf{Solvanta Engineering Solutions} | DPyPQ}
\rhead{\small Documento Oficial TE-01 -- Tabla de Energía y Servicios}
\cfoot{\small Página \thepage\ de \pageref{LastPage}}
\renewcommand{\headrulewidth}{0.4pt}

\begin{document}

\begin{center}
{\large \textbf{SOLVANTA ENGINEERING SOLUTIONS S.A. DE C.V.}}\\[0.15cm]
{\Large \textbf{\color{navyblue} DOCUMENTO TE-01: TABLA DE CONSUMO DE ENERGÍA Y SERVICIOS AUXILIARES}}\\[0.1cm]
{\small Planta de Producción de e-Metanol ($30\,000\text{ t/año}$) | Caso Base Definitivo DWSIM v1 ($513.0\text{ kmol/h }\text{H}_2\text{O}$ / $121.0\text{ kmol/h }\text{CO}_2$)}
\end{center}

\vspace{0.2cm}
\setlength{\tabcolsep}{4pt}
\begin{xltabular}{\textwidth}{l p{2.6cm} p{4.2cm} c r r r >{\raggedright\arraybackslash}X}
\caption{\textbf{TE-01. Resumen de consumos eléctricos y cargas térmicas de planta (DWSIM v1 exacto).}} \label{tab:te01} \\
\toprule
\textbf{Tag} & \textbf{Servicio} & \textbf{Descripción} & \textbf{$T_{\text{in}} \to T_{\text{out}}$} & \textbf{$Q$ (kW)} & \textbf{$W$ (kW)} & \textbf{MWh/año} & \textbf{Observación Técnica / Suministro} \\
\midrule
\endfirsthead
\toprule
\textbf{Tag} & \textbf{Servicio} & \textbf{Descripción} & \textbf{$T_{\text{in}} \to T_{\text{out}}$} & \textbf{$Q$ (kW)} & \textbf{$W$ (kW)} & \textbf{MWh/año} & \textbf{Observación Técnica / Suministro} \\
\midrule
\endhead
\bottomrule
\multicolumn{8}{r}{\footnotesize\textit{Continúa en la siguiente página\dots}} \\
\endfoot
\bottomrule
\endlastfoot
\texttt{EL-201} & Electricidad & Stack electrólisis SOEC & 800 -> 800 & --- & +38\,067.30 & 304\,538.4 & Escenario D Óptimo (800°C, 0.480 A/cm², 1.183 V, 38.07 MWe, 32.92 kWh/kg H2). \\
\texttt{E-202} & Calef. eléctr. & Sobrecalentador vapor SOEC & 220 -> 800 & +5\,248.42 & --- & 41\,987.4 & Calentamiento vapor a 20.0 bar (5.25 MWth). \\
\texttt{E-203} & Agua enfriam. & Condensador efluente SOEC & 800 -> 38 & -8\,962.22 & --- & 71\,697.8 & Enfriamiento y condensación H2O catódica. \\
\texttt{E-301} & Calefacción & Precalentador syngas fresco & 33 -> 72 & +250.00 & --- & 2\,000.0 & Acondicionamiento térmico previo a C-301. \\
\texttt{C-301} & Electricidad & Compresor syngas multietapa & 72 -> 271 & --- & +1\,316.67 & 10\,533.4 & Presurización de 19.4 a 75.0 bar (1.32 MWe). \\
\texttt{R-401-S} & Gen. vapor & Carcasa caldera BWR & 40 -> 220 & -12\,335.36 & --- & 98\,682.9 & Genera vapor saturado 20 bar para SOEC. \\
\texttt{R-401} & Reac. exotérmica & Reactor tubular Lurgi BWR & 245 -> 245 & -7\,924.76 & --- & 63\,398.1 & Calor neto de reacción química desprendido. \\
\texttt{E-502A/B} & Integr. térm. & Economizador FEHE & 245 -> 225 & +1\,235.43 & --- & 9\,883.4 & Precalienta reciclo a 60°C con efluente. \\
\texttt{E-501} & Agua enfriam. & Condensador síntesis & 225 -> 38 & -13\,742.75 & --- & 109\,942.0 & Condensa mezcla metanol/agua previo a V-501. \\
\texttt{C-501} & Electricidad & Soplador de reciclo & 60 -> 60 & --- & +150.00 & 1\,200.0 & Circulación de gas bucle (delta P = 1.5 bar). \\
\texttt{X-501} & Gen. eléctr. & Turboexpansor de purga & 60 -> -139 & --- & -437.79 & -3\,502.3 & Expansión 75 -> 1.1 bar (recuperación energía). \\
\texttt{E-601} & Agua enfriam. & Condensador domo T-601 & 68.5 -> 68.5 & -3\,077.25 & --- & 24\,618.0 & Condensación a reflujo (R/D = 1.50). \\
\texttt{E-602} & Vapor baja & Rehervidor torre T-601 & 104 -> 104 & +3\,266.68 & --- & 26\,133.4 & Ebullición fondos columna fraccionadora. \\
\texttt{E-603} & Agua helada & Subenfriador producto & 68.5 -> 35 & -111.11 & --- & 888.9 & Acondicionamiento final a tanque TK-601. \\
\texttt{P-801} & Electricidad & Bomba retorno ZLD & 59 -> 59 & --- & +5.51 & 44.1 & Presurización agua tratada a 20.0 bar. \\
\midrule
\multicolumn{4}{l}{\textbf{Demanda Eléctrica Neta de Planta}} & --- & $\mathbf{+39\,101.69}$ & $\mathbf{312\,813.5}$ & Potencia neta de proceso ($39.10\text{ MWe}$). \\
\bottomrule
\end{xltabular}

\vspace{0.3cm}
\noindent\textbf{Balance Energético Consolidado:}
\begin{itemize}
    \item \textbf{Potencia Eléctrica Total Requerida:} $\mathbf{39\,101.69\text{ kW}}$ ($\mathbf{39.10\text{ MWe}}$), absorbiendo la electrólisis SOEC el $98.1\%$ de la potencia demandada.
    \item \textbf{Generación Eléctrica por Recuperación en Turboexpansor \texttt{X-501}:} $\mathbf{437.79\text{ kW}}$ continuos inyectados a la red interna de media tensión.
    \item \textbf{Recuperación de Calor de Reacción Exotérmica:} $\mathbf{12\,335.36\text{ kW}}$ transferidos a la generación de vapor saturado a $20.0\text{ bar}$ en la carcasa caldera \texttt{R-401-S}.
\end{itemize}

\end{document}
